\vspace{-2mm}
\section{Related Works}
\vspace{-1mm}
\subsection{Video World Models}
% \vspace{-2mm}
\label{sec:related_world_modeling}
World modeling aims to construct an evolving environment to simulate the real world. Existing works approach this problem from different perspectives. \cite{marble} builds explicit geometric-consistent 3D representations, while JEPA \cite{vjepa, vjepa2} learns abstract state transitions in latent space. With the rapid progress of video generation models~\cite{svd,animatediff,wan,hunyuanvideo}, Generative Video World Models \cite{genie1,genie3,waymo} have become a dominant paradigm as they achieve more scalable and realistic world modeling. Typically, a video world model predicts future frames conditioned on historical context and control signals. First, to leverage historical content, early methods \cite{svd,animatediff,wan,hunyuanvideo,history} concatenate anchor frames at the sequence start. CausVID \cite{CausVID} pioneeringly distills fixed-size self-attention into causal attention, enabling autoregressive next-frame prediction with KV cache \cite{selfforcing, selfforcing++,longlive,rollingforcing,deepforcing,blockvid, contextasmemory,genie1,genie3,liveavatar} to store arbitrary history tokens. Next, to support controllable exploration, recent works \cite{hunyuangamecraft,yan,pan,relic,matrixgame, cameractrl1,cameractrl2} incorporate camera trajectory embeddings. Combined with cached history, these methods allow revisiting previously observed regions. However, the cached frame tokens merely capture 2D visual snapshots of past regions at the exact time they were observed. This inherently relies on a static-world assumption, completely failing to maintain out-of-sight dynamics.

\vspace{-1.5mm}
\subsection{Explicit Spatial Memory}
To achieve precise camera control and maintain long-term geometry during exploration, recent methods \cite{spatia, vmem,hunyuanworld,hunyuanworld15,worldplay,spatialmem,epic,anchorweave} have been developed conditioned on explicit 3D spatial memory. These approaches maintain an explicitly reconstructed spatial representation as a form of global memory---such as a point cloud combined with camera parameters predicted by feed-forward estimators \cite{wang2025vggt,wang2024dust3r,wang2025continuous,depthanything3,keetha2026mapanything,stream3r2025}. By injecting this structural representation into the video generation model, they ensure rigorous geometric consistency.
Despite these advancements in spatial tracking, the temporal dimension of historical locations remains neglected. The registered 3D representations store merely the static 3D spatial structure of the scene at the moment of observation. Consequently, they still cannot capture the \textit{out-of-sight dynamics} within previously visited areas.

\vspace{-2.5mm}
\subsection{Out-of-sight Dynamics}
\vspace{-1mm}
An ideal world model should capture the continuous evolution of the entire explored space, forming a unified 4D spatio-temporal field where scene states evolve consistently over time. This poses a fundamental challenge to existing video world models. Current approaches~\cite{selfforcing, selfforcing++,longlive,rollingforcing,deepforcing,blockvid, contextasmemory,genie1,genie3,liveavatar, spatia, vmem,hunyuanworld,hunyuanworld15,worldplay,spatialmem,epic,anchorweave} only update states within the camera’s visible region, while content outside the current view is merely stored as historical observations in memory (e.g., KV cache or spatial memory), remaining frozen at their last observed timestamps. In this work, we formalize this overlooked problem within the video world model paradigm. Building upon explicit spatial memory, we propose a decoupled framework to continuously maintain and evolve this out-of-view representation, thereby bridging the gap between static scene memorization and true 4D dynamic world simulation.